% Conversão de chapter-06-text-h-0106; origem RELATORIO_ATUAL.odt.
\hypertarget{__RefHeading___Toc4829_3118176633}{}
\refstepcounter{section}
\section*{5.1 Introdução}
\label{ages-iv:text-h-0106}
\addcontentsline{toc}{section}{5.1 Introdução}

% SOURCE text-p-1325
O consumo de música está cada vez mais integrado às plataformas digitais, que permitem descobrir artistas, acompanhar o que outras pessoas estão ouvindo e compartilhar conteúdos. Apesar disso, essas experiências normalmente não registram as histórias, emoções e momentos pessoais que fazem com que determinadas músicas adquiram significado para cada indivíduo. Nesse contexto, o Recorda surge a partir da lacuna entre o consumo musical e a memória afetiva, propondo uma forma de conectar músicas a experiências pessoais por meio de relatos, fotos e vídeos.\par

% SOURCE text-p-1326
O projeto consiste no desenvolvimento de um aplicativo mobile para iOS e Android voltado ao registro, organização e compartilhamento de memórias musicais, denominadas Recordas. A plataforma busca combinar características de diário digital e rede social, permitindo associar experiências a músicas, artistas e álbuns, além de oferecer recursos como integração com serviços de streaming, feeds de conteúdo, interações sociais, recomendações por afinidade musical e estatísticas.\originalspace{}\par

% SOURCE text-p-1327
Além do registro das memórias, o Recorda contempla funcionalidades voltadas à descoberta e à interação entre usuários, como amizades, recomendações, retrospectivas e compartilhamento de conteúdos. A solução também prevê um painel administrativo para gestão e moderação da plataforma e suporte aos idiomas português, inglês e espanhol.\par

% FIGURE draw-image-0024
\memorialfigure{imagens/original/asset-0011.jpg}{\textbf{Figura 21 –} Foto do Time - AGES IV}{ages-iv:draw-image-0024}{Fonte: pendente de identificação}{0.92}
